\subsection{Случайные величины.}
\subsubsection{Определение случайной величины, измеримые и борелевские функции.}

Зафиксируем измеримое пространство $(\Omega, \mathcal{F}, \P)$ и отображение $f: \omega \xrightarrow{f}\R$.

\begin{definition}
    Отображение $\xi = f(\omega)$ или просто $\xi(\omega)$ будем называть случайной величиной при cледующего условия:
    \begin{gather*}
        \forall x \ D_x = \{\omega : \xi(\omega) < x \} \equiv \{ \omega : \xi(\omega) \in (-\infty, x) \} \in \mathcal{F}.
    \end{gather*}
    Другими словами, случайная величина есть измеримое отображение из $\Omega$ в $\R$.
\end{definition}
\begin{remark}
    Множество $\overline{D}_x = \{\omega : \xi(\omega) \in [x, + \infty) \} \in \mathcal{F}$.

    \noindent Пусть $x < y$. Тогда $D_y\overline{D}_x = \{\omega : \xi(\omega) \in [x, y) \} \in \mathcal{F}$.

    \noindent Следовательно, $\{\omega : \xi(\omega) \in B \in \mathcal{B} \} \in \mathcal{F}$.
\end{remark}
\begin{definition}
    Сигма-алгеброй, порожденной случайной величиной $\xi$ будем называть следующую сигма-алгебру $\sigma(\xi) = \sigma_{min}(\{D_x \}_{x \in \R})$.
\end{definition}

\begin{definition}
    Функцию $\xi$ будем называть измеримой, если $\forall B \in \mathcal{B} \hookrightarrow \xi^{-1}(B) \in \mathcal{F}$.
\end{definition}
\begin{example}
    Рассмотрим $\xi(\omega) \equiv C.$ Тогда $D_x = \{\varnothing, \Omega\} \equiv \sigma(\xi).$
\end{example}

\begin{property}
    Перечислим свойства случайных величин:
    \begin{enumerate}
        \item $\eta = \phi(\xi)$ тоже является случайной величиной, если $\phi : \R \rightarrow \R$ и $\phi^{-1}(B) \in \mathcal{B}$.
        \item Если $\xi$~--- случайная величина, то $a\xi + b, |\xi|, \xi^{\alpha} (\alpha > 0), e^{\xi}$~--- случайные величины. 
    \end{enumerate}
\end{property}
\subsubsection{Измеримость \texorpdfstring{$\sup$}{sup} и \texorpdfstring{$\inf$}{inf}, суммы и произведения случайных величин.}
\begin{property}
    Пусть $\xi_1(x), \xi_2(\omega)$ ~---~ случайные величины. Покажем, что $\xi_1(\omega ) + \xi_2(\omega)$~--- случайная величина.
\end{property}
\begin{proof}
    Хотим показать, что $\{\omega : \xi_1(\omega) + \xi_2(\omega) < x \} \in \mathcal{F}$.

    \noindent Рассмотрим множество $\{ \omega : \xi_1(\omega) < \xi_2(\omega) \} = \bigcup\limits_{k} \{ \omega : \xi_1(\omega) < r_k \} \cap \{\omega: \xi_2(\omega) \ge r_k \}$, где $r_k = \Q$.

    \noindent Оба множества в объединении измеримы. Тогда верно следующее: 
    \begin{gather*}
        \{\omega : \xi_1(\omega) + \xi_2(\omega) < x \} = \{ \omega : \xi_1(\omega) < x - \xi_2(\omega) \} \in \mathcal{F}.
    \end{gather*}
\end{proof}

\begin{property}
    Умножать случайные величины тоже можно, так как $\xi_1 + \xi_2 = \frac{1}{4}((\xi_1 + \xi_2)^2 + (\xi_1 - \xi_2)^2)$.
\end{property}

\begin{property}
    Поймем можно ли делить. $\dfrac{\xi_1}{\xi_2} = \xi_1 \cdot \frac{1}{\xi_2}$, если $\xi_2(\omega) \neq 0, \forall \omega$.
\end{property}

\subsubsection{Критерий существований функциональной зависимости случайных величин.}
\begin{theorem}
    Пусть заданы две случайные величины $\xi$ и $\eta$.
    Тогда следующие условия эквивалентны:
    \begin{enumerate}
        \item $\sigma(\xi) \subset \sigma(\eta)$;
        \item Существует такая борелевская $g$, что $\xi = g(\eta)$.
    \end{enumerate}
\end{theorem}
\begin{proof}
    Импликация $2) \Ra 1)$ очевидна. \\
    Покажем справедливость $1) \Ra 2)$. 
    Определим последовательность $A_{k, n}$ как
    \[
        A_{k, n} = \left\{\xi \in \left[\dfrac{k}{2^n}, \dfrac{k + 1 }{2^{n}}\right]\right\}.
    \]
    В силу измеримости $\xi$ относительно $\eta$ справедливо, что $A_{k, n} \in \sigma(\eta)$.
    Пусть $B_{k, n} = \{\eta(w) \mid w \in A_{k, n}\}$ и функция $g_n(x) = \sum\limits_{k, n} \dfrac{k}{n}\Ind_{B_{k, n}}$. Заметим, что $g_n$ -- возрастающая последовательность простых функций и $|g_n(\eta) - \xi| \ra 0, n \ra \infty$, что и показывает существование борелевской $g$ такой, что $g(\eta) = \xi$.
\end{proof}